%% kapa: exam=1 round=20 year=2009 period=2 form=A subject=appraisal_law subject_ko=부동산관계법규 seq=7 number=47
%% source: Q-Net 감정평가사 기출문제 (cst00309) artlSeq=5024178 file=extracted/1008/5024178/2009년 제20회 감정평가사 1차시험 기출문제/2교시.hwp
%% conversion: kordoc 4.24.51 -> tools/gen1.py (auto); review=proofread-sonnet; flags=none
\begin{question}{47}
국토의 계획 및 이용에 관한 법령상 선매에 관한 설명으로 옳지 \uline{않은} 것은?
\begin{choices}
  \item 시장ㆍ군수ㆍ구청장은 토지거래계약에 관한 허가 신청이 있는 경우 공익사업용 토지에 대해서는 선매자를 지정하여 당해 토지를 협의 매수하게 할 수 있다.
  \item 국가는 토지거래계약에 관한 허가 신청이 있는 경우 공익사업용 토지에 대한 선매자가 될 수 있다.
  \item 선매자가 토지를 매수하는 경우 가격은 부동산 가격 공시 및 감정 평가에 관한 법률에 따른 공시지가로 하여야 한다.
  \item 시장ㆍ군수ㆍ구청장은 공익사업용 토지에 대하여 토지거래계약 허가 신청이 있는 날로부터 1월 이내에 선매자를 지정하여 토지소유자에게 이를 통지하여야 한다.
  \item 시장ㆍ군수ㆍ구청장은 선매 협의가 이루어지지 아니한 때에는 지체 없이 허가 또는 불허가의 여부를 결정하여 이를 통보하여야 한다.
\end{choices}
\end{question}
\answer{3}
